\documentclass[11pt,a4paper]{article}

% --- Packages ---
\usepackage[utf8]{inputenc}
\usepackage[T1]{fontenc}
\usepackage{geometry}
\geometry{margin=2cm}
\usepackage{array}
\usepackage{fancyhdr}
\usepackage{titlesec}
\usepackage{tcolorbox}
\tcbuselibrary{listings,skins,breakable}
\usepackage{amsmath,amssymb}
\usepackage{float}
\usepackage{xurl}
\usepackage[hidelinks]{hyperref}
\usepackage{graphicx}
\usepackage{subcaption}
\usepackage[spanish]{babel}

\lstdefinestyle{pythonstyle}{
  language=Python,
  backgroundcolor=\color{gray!5},
  basicstyle=\ttfamily\small,
  keywordstyle=\color{blue}\bfseries,
  commentstyle=\color{gray}\itshape,
  stringstyle=\color{teal},
  numberstyle=\tiny\color{gray},
  numbers=left,
  stepnumber=1,
  numbersep=8pt,
  showspaces=false,
  showstringspaces=false,
  showtabs=false,
  tabsize=4,
  breaklines=true,
  breakatwhitespace=true,
  frame=single,
  rulecolor=\color{gray!40},
  captionpos=b,
  keepspaces=true,
  morekeywords={ceil,cbrt,round,split},
}


% --- Formatting ---
\pagestyle{fancy}
\fancyhf{}
\setlength{\headheight}{30pt}
\setlength{\footskip}{20pt}
\setlength{\parindent}{0pt}
\setlength{\parskip}{0.7em}

% Left header (university and exam info)
\rhead{%
   \small
   Universidad Jorge Tadeo Lozano\\
   Estructuras de Datos - Grupo 2\\
   Prueba de Programación II 
}

% Right header (date, professor, monitor)
\lhead{%
   \small
   \textbf{Fecha}: \today \\
   \textbf{Profesor:} Alejandro Franco \\
   \textbf{Monitor:} Ludwig Alvarado
}

\cfoot{\thepage}

\titleformat{\section}{\large\bfseries}{\thesection.}{0.5em}{}

% --- Environment for samples ---
\newenvironment{sample}[2]{
  \begin{tcolorbox}[title=Prueba \##1, colback=white, colframe=black, sharp corners,
    fonttitle=\bfseries, breakable]
    \begin{tabular}{|p{0.45\linewidth}|p{0.45\linewidth}|}
      \hline
      \textbf{Entrada} & \textbf{Salida} \\
      \hline
      }{\\ \hline
    \end{tabular}
  \end{tcolorbox}
}

% --- Document ---
\begin{document}

\begin{center}
  {\LARGE \bf Vibing Insecurities}\\[1ex]
  Autor: Ludwig Alvarado
\end{center}

La Unidad Técnica de Agentes, Detectives y Operaciones (UTADEO) es una organización clandestina que se dedica al espionaje académico y a la manipulación de registros de calificaciones. Muchos estudiantes recurren a sus servicios para alterar sus notas, mientras que la organización identifica posibles clientes analizando sus inseguridades durante los exámenes. A este proceso lo denominan \textit{vibing insecurities}.

Además, UTADEO se aprovecha del auge del \textit{vibe coding}, donde múltiples instituciones han delegado el desarrollo de sus sistemas a modelos de lenguaje y programadores poco rigurosos. Estas vulnerabilidades permiten que el equipo de inyección SQL modifique la base de datos académica sin ser detectado.

Tú haces parte de este equipo. Tu misión es implementar un sistema que procese una serie de operaciones sobre los registros de estudiantes y sus calificaciones.

Cada estudiante tiene:
\begin{itemize}
    \item Un nombre único
    \item Una lista de calificaciones
\end{itemize}

El puntaje final de un estudiante se define como el promedio de sus calificaciones.

\subsection*{Entrada del programa}

La primera línea contiene un entero $N$, el número inicial de estudiantes.

Las siguientes $N$ líneas contienen la información de cada estudiante en el formato:

\texttt{name k score\_1 score\_2 ... score\_k}

Donde:
\begin{itemize}
    \item \texttt{name} es una cadena en minúsculas
    \item $k$ es la cantidad de calificaciones
\end{itemize}

Luego se proporciona un entero $Q$ — el número de consultas.

Las siguientes $Q$ líneas contienen una de las siguientes operaciones:

\begin{itemize}
    \item \texttt{ADD name k score\_1 ... score\_k}
    \item \texttt{UPDATE name k score\_1 ... score\_k}
    \item \texttt{APPEND name score}
    \item \texttt{TOP k}
    \item \texttt{AVERAGE name}
    \item \texttt{GLOBAL\_AVERAGE}
\end{itemize}

\subsection*{Descripción de las operaciones}

A continuación se describe el comportamiento exacto de cada operación:

\begin{itemize}

\item \texttt{ADD name k score\_1 ... score\_k}

Agrega un nuevo estudiante con nombre \texttt{name} y una lista de $k$ calificaciones.

\begin{itemize}
    \item Si el estudiante no existe, se inserta en el sistema con sus calificaciones.
    \item Si el estudiante ya existe, la operación se ignora y no se realiza ningún cambio.
\end{itemize}

\vspace{0.3cm}

\item \texttt{UPDATE name k score\_1 ... score\_k}

Reemplaza completamente la lista de calificaciones de un estudiante existente.

\begin{itemize}
    \item Si el estudiante existe, su lista de calificaciones anterior se elimina y se reemplaza por la nueva.
    \item Si el estudiante no existe, la operación se ignora.
\end{itemize}

\vspace{0.3cm}

\item \texttt{APPEND name score}

Agrega una nueva calificación al final de la lista de un estudiante.

\begin{itemize}
    \item Si el estudiante existe, la calificación se añade a su lista.
    \item Si el estudiante no existe, la operación se ignora.
\end{itemize}

\vspace{0.3cm}

\item \texttt{TOP k}

Consulta los $k$ estudiantes con mejor desempeño según los siguientes criterios:

\begin{enumerate}
    \item Mayor promedio de calificaciones
\end{enumerate}

\begin{itemize}
    \item Se deben imprimir exactamente $k$ estudiantes o todos los existentes si hay menos de $k$.
    \item Cada estudiante se imprime en una línea en el formato:
    
    \texttt{name average}
    
    \item El promedio debe mostrarse con exactamente 2 decimales.
\end{itemize}

\vspace{0.3cm}

\item \texttt{AVERAGE name}

Consulta el promedio de un estudiante específico.

\begin{itemize}
    \item Si el estudiante existe, se imprime su promedio con exactamente 2 decimales.
    \item Si el estudiante no existe, se imprime:
    
    \texttt{NA}
\end{itemize}

\vspace{0.3cm}

\item \texttt{GLOBAL\_AVERAGE}

Calcula el promedio global considerando todas las calificaciones de todos los estudiantes.

\begin{itemize}
    \item Se deben sumar todas las calificaciones y dividir entre el total de calificaciones.
    \item El resultado se imprime con exactamente 2 decimales.
    \item Se garantiza que siempre existe al menos una calificación en el sistema al momento de esta consulta.
\end{itemize}

\end{itemize}

\subsection*{Salida del programa}

Para cada consulta que genere salida, imprime el resultado en una nueva línea.

\begin{itemize}
    \item \texttt{TOP k}: imprime los $k$ mejores estudiantes, cada uno en una línea como:
    
    \texttt{name average}
    
    \item \texttt{AVERAGE name}: imprime el promedio con 2 decimales, o \texttt{NA} si no existe

    \item \texttt{GLOBAL\_AVERAGE}: imprime el promedio global con 2 decimales
\end{itemize}

\textbf{Nota:}
\begin{itemize}
    \item \texttt{ADD}, \texttt{UPDATE} y \texttt{APPEND} no generan salida
    \item Las salidas deben imprimirse en el orden de las consultas
\end{itemize}

\subsection*{Restricciones}

\begin{itemize}
    \item $1 \leq N, Q \leq 10^5$
    \item $1 \leq k \leq 100$
    \item $0 \leq score \leq 10^6$
    \item Longitud del nombre $\leq 20$
\end{itemize}

\subsection*{Caso de prueba de ejemplo}

\begin{sample}{1}{}
\texttt{2}

\texttt{alice 3 90 80 100}

\texttt{bob 2 70 80}

\texttt{7}

\texttt{TOP 2}

\texttt{APPEND bob 100}

\texttt{TOP 2}

\texttt{AVERAGE alice}

\texttt{GLOBAL\_AVERAGE}

\texttt{UPDATE alice 2 50 50}

\texttt{TOP 2}
&
\texttt{alice 90.00}

\texttt{bob 75.00}

\texttt{bob 83.33}

\texttt{alice 90.00}

\texttt{90.00}

\texttt{86.67}

\texttt{bob 83.33}

\texttt{alice 50.00}
\end{sample}

\subsection*{Notas}

\begin{itemize}
    \item El promedio se calcula como $\frac{\sum scores}{|scores|}$
    \item Todas las salidas numéricas deben imprimirse con exactamente 2 decimales
    \item Si una operación hace referencia a un estudiante inexistente, debe ignorarse (excepto en \texttt{AVERAGE}, donde se imprime \texttt{NA})
    \item Se recomienda evitar recalcular promedios desde cero en cada consulta por eficiencia
\end{itemize}

\subsection*{Aspectos Importantes}

\begin{enumerate}
  \item La prueba se resolverá en \textbf{2 momentos}:
        \begin{itemize}
            \item Desarrollo del algoritmo en papel \textbf{(3.0 puntos)}.
            \item Implementación en el editor de código \textbf{(2.0 puntos)}.
        \end{itemize}
  \item Para leer los datos puede utilizar la función \lstinline[style=pythonstyle]|split()| de la clase \lstinline[style=pythonstyle]|str| y de ahí indexar los datos en un diccionario con llave de tipo \lstinline[style=pythonstyle]|str| y valor de tipo \lstinline[style=pythonstyle]|list|. Por ejemplo, para procesar la línea \texttt{alice 3 90 80 100} esto se puede guardar dentro de un diccionario como: 
        \begin{lstlisting}[style=pythonstyle]
>>> entrada = input().split()
>>> diccionario = {}
>>> diccionario[entrada[0]] = list(map(int, entrada[1::]))
>>> print(diccionario)
{"alice": [3, 90, 80, 100]}
        \end{lstlisting}

  \item Recuerda que la función \lstinline[style=pythonstyle]|round()| ayuda a redondear números decimales. Por ejemplo, el número \texttt{3.141592653589793} al aplicar \lstinline[style=pythonstyle]|round(3.141592653589793, 2)| esto retorna en \texttt{3.14}

  \item Este ejercicio evalúa los temas vistos hasta el \textbf{segundo corte}.

  \item No es necesario usar estructuras de datos avanzadas ni importar librerías externas, \textbf{ante la duda pregunte al profesor o monitore, de lo contrario, su nota se verá afectada.}

  \item No es necesario validar errores de entrada.

  \item La única conversión de tipo de datos debe ser al momento de la lectura de los mismos. No es necesario volver a convertir el tipo de dato de las variables.

  \item El formato de salida debe coincidir exactamente con el especificado. Cualquier diferencia será considerada incorrecta.
\end{enumerate}





\end{document}
